\subsubsection{Multidimensional Dynamic Attention (MDA)}

\citet{almaghrabi2024multidimensional} construct two synthetic forecasting
systems to evaluate whether multidimensional dynamic attention recovers the
lagged variables used by the target equation. The two systems deliberately
pose different identification problems.

\paragraph{Toy1: controlled relevance with independent inputs.}
Toy1 contains six predictor series sampled independently from uniform ranges:
$x_1\in[2,5]$, $x_2\in[10,30]$, $x_3\in[5,10]$, $x_4\in[30,70]$,
$x_5\in[100,200]$, and $x_6\in[65,85]$. From these predictors, the paper
generates 20,000 target samples with

\begin{equation}
    y_t
    = x_{1,t-1}x_{2,t-2}
    + x_{3,t-5}
    + \sum_{\ell\in\{1,4,5,7,8\}}x_{4,t-\ell}
    + \varepsilon_t,
\end{equation}

where $\varepsilon_t$ is Gaussian white noise with mean zero and standard
deviation 0.1 (paper Eq.~15). The exact relevant locations are $x_1$ at lag 1,
$x_2$ at lag 2, $x_3$ at lag 5, and $x_4$ at lags 1, 4, 5, 7, and 8. The
product $x_{1,t-1}x_{2,t-2}$ also tests whether the explanation identifies both
participants in a nonlinear interaction. Variables $x_5$ and $x_6$ are
deliberately included as irrelevant predictors. Thus, the reported $V=7$
counts six predictors plus the target series $y$. Historical values of $y$ can
therefore be supplied to the forecaster even though they are not explicit terms
of paper Eq.~15; the experiment uses an input length of 20 and predicts eight
future target values.

\paragraph{Toy2: relevance within chaotic time-dependent inputs.}
Toy2 replaces the independently sampled predictors with four independently
generated Mackey--Glass series. These inputs have their own nonlinear delayed
temporal dynamics. The paper uses the parameter triples $(\beta,\gamma,\tau)$
$x_1=(0.1,0.2,17)$, $x_2=(0.2,0.3,17)$, $x_3=(0.1,0.2,20)$, and
$x_4=(0.2,0.4,5)$, and then generates the target through

\begin{equation}
    y_t
    = x_{1,t-2}x_{2,t-6}
    + x_{3,t-4}
    + x_{4,t-1}x_{4,t-2}
    + \varepsilon_t,
\end{equation}

which is paper Eq.~16. Its binary reference marks $x_1$ at lag 2, $x_2$ at
lag 6, $x_3$ at lag 4, and $x_4$ at lags 1 and 2. All four predictor series
participate in the target equation; unlike Toy1, Toy2 has no deliberately
irrelevant predictor among $x_1,\ldots,x_4$. The two products introduce two
interactions, including an interaction between two lags of the same series
$x_4$. Here $V=5$ counts the four Mackey--Glass predictors plus $y$; the longer
input window contains 60 values and the model predicts 20 future target values.
In the discussion of Fig.~8, the authors call the target the fifth series
$x_5$ and state that its lagged values are also model inputs. Those target lags
may be predictively useful because of temporal dependence, but they are not
explicit driver terms in the structural reference obtained from Eq.~16.

The central distinction is therefore that Toy1 tests sparse lag recovery and
suppression of irrelevant variables when the predictors themselves are
independently sampled, whereas Toy2 tests lag recovery in predictor series that
already contain chaotic delayed dynamics, using a longer input window and
forecast horizon. Toy1 contains eight relevant variable--lag locations and two
irrelevant variables; Toy2 contains five relevant variable--lag locations and
uses every predictor, although all remaining variable--lag locations are still
irrelevant.

For both equations, $t$ denotes the target time. The latest available input for
the one-step target $y_t$ is at $t-1$: this is $x_{1,t-1}$ and $x_{4,t-1}$ in
Toy1, and $x_{4,t-1}$ in Toy2. With forecast origin $\tau=t-1$, those samples
have relative position 0, while $t-2$ has relative position $-1$. For a later
forecast step, the same generator lags shift by one position in the fixed input
window for every additional step into the horizon.

Both equations yield a binary structural reference over the variable--lag
grid: every location occurring in the applicable target equation is relevant,
and every other location is irrelevant. The reference records participation,
including both inputs of each multiplicative term, but does not assign a signed
local contribution magnitude. SHAPformer's reference instead assigns a signed
local contribution to every component and forecast step.

\paragraph{Original dataset-configuration table.}
Table~\ref{tab:mda-dataset-settings-exact} reproduces Table~4 of
\citet{almaghrabi2024multidimensional}. The two synthetic datasets are
$Toy_1$ and $Toy_2$; NSW, QLD, and Pollution are the three real datasets listed
in the same source table. The generating mechanisms for $Toy_1$ and $Toy_2$
are supplied separately by Eqs.~15--16 of the paper, while this table fixes the
input length, output horizon, and number of variables used in the experiments.

\begin{table}[htbp]
\centering
\small
\caption{Exact transcription of Table~4 in
\citet{almaghrabi2024multidimensional}: summary of the length of input data
($X$), the length of predicted output data ($Y$), and the number of variables
($V$) in the input sequence for each dataset.}
\label{tab:mda-dataset-settings-exact}
\begin{tabular}{@{}lccccc@{}}
\toprule
Dataset & $Toy_1$ & $Toy_2$ & NSW & QLD & Pollution \\
\midrule
$\operatorname{len}(X)$ & 20 & 60 & 27 & 54 & 48 \\
$\operatorname{len}(Y)$ & 8  & 20 & 27 & 27 & 24 \\
$V$                     & 7  & 5  & 14 & 14 & 8  \\
\bottomrule
\end{tabular}
\end{table}

\paragraph{How MDA is evaluated.}
The final feature--time score for forecast step $h$ is the elementwise product
of variable and temporal attention scores,
$\alpha^{\mathrm{c}}_{h,j,\ell}
=\alpha^{\mathrm{va}}_{h,j,\ell}\alpha^{\mathrm{ta}}_{h,j,\ell}$
(Eq.~12 of \citet{almaghrabi2024multidimensional}). The authors plot these
scores for Toy1 at $h=1,2,4$ (Fig.~7) and Toy2 at $h=2,20$ (Fig.~8), alongside
RF, MAFS, and TFT. They inspect whether high weights coincide with known
variable--lag locations, whether the deliberately irrelevant Toy1 variables
$x_5,x_6$ receive little weight, and whether the location changes with $h$
as expected from the equation. This is a visual comparison of continuous
scores with known structural support; precision, recall, F1, correlation, and
other numerical attention-to-ground-truth errors lie outside the reported
evaluation. The
multiplicative Toy1 term establishes that $x_{1,t-1}$ and $x_{2,t-2}$ both
belong to the reference support; the plots treat them as two relevant locations
and leave recovery of their interaction effect unquantified.

Separately, MAE and RMSE compare predicted and observed target values
(Eqs.~17--18; Table~5); attention agreement is assessed through the maps above.
The architecture
ablation removes variable attention, temporal attention, or dynamic
representation learning and compares forecasting MAE/RMSE (Table~6).
In the official implementation, these metrics are computed after selecting
the last forecast step from each predicted sequence; the explanation evaluation
therefore remains map-based.\footnote{Official MDA code:
\url{https://github.com/sarah-almaghrabi/MDA-Multidimensional-Dynamic-Attention};
the horizon selection is in \texttt{main.py}.}

